\documentclass[11pt]{article}
\usepackage{fontspec}
\setmainfont{Times New Roman}
\setsansfont{Arial}
\setmonofont{Consolas}
\usepackage{amsmath,amssymb}
\usepackage{geometry}
\usepackage{microtype}
\geometry{a4paper,margin=3cm}
\setlength{\parindent}{1.25em}
\setlength{\parskip}{0pt}
\emergencystretch=10em
\allowdisplaybreaks
\providecommand{\Kfield}{\mathfrak{K}}
\providecommand{\Ssys}{\mathfrak{S}}
\providecommand{\Mfield}{\mathfrak{M}}
\providecommand{\tq}{\tau}
\providecommand{\ds}{\displaystyle}

\begin{document}

\begin{center}
{\Large\bfseries 5. Поля раціональних функцій}\par
\vspace{1.2em}
J. Ber. d. DMV 22 (1913), S. 316--319
\end{center}

\vspace{2em}

Наведені нижче питання первісно походять із розмов з Е. Фішером. Деякі з них, до речі, для спеціального випадку однієї невизначеної -- і навіть за загальніших припущень щодо поля коефіцієнтів -- уже поставив і розв'язав Е. Штайніц.\footnote{E. Steinitz: Algebraische Theorie der Körper (особливо \S\ 24). Crelles Journal 137. 1910.}

Під ``полем раціональних функцій'' я розумію поле, елементи якого є раціональними функціями від $n$ невизначених із коефіцієнтами з наперед заданого числового поля, яке, зокрема, може містити і всі комплексні числа. Прикладами таких полів є поле симетричних функцій від $n$ величин або, загальніше, сукупність раціональних функцій від $n$ невизначених, що допускають перестановки певної групи (Лагранжеві області родів); далі -- поле інваріантів.

Між двома першими названими полями і останнім є характерна відмінність: перші містять $n$ алгебраїчно незалежних функцій, елементарні симетричні функції; тоді як кількість алгебраїчно незалежних інваріантів завжди менша за кількість невизначених. Тому важливо, що загалом такі поля другого типу можна взаємно однозначно зіставляти з полями першого типу, замінюючи деякі з невизначених числами. Отже, далі я обмежуватимусь полями першого типу, тобто полями з $n$ алгебраїчно незалежними функціями; завдяки цьому зіставленню результати тоді мають загальну силу.

Питання групуються навколо трьох базових понять: раціональна база, мінімальна база та інтегральна база.

1. Під ``раціональною базою'' я розумію скінченну кількість функцій поля таку, що кожну функцію поля можна подати як раціональну комбінацію цієї скінченної кількості функцій з коефіцієнтами із заданої області коефіцієнтів. Простими міркуваннями -- які у випадку однієї невизначеної містяться також у Е. Штайніца, там само -- доводять існування раціональної бази для кожного поля раціональних функцій. Наприклад, за теоремою Лагранжа елементарні симетричні функції і одна функція, що належить до групи, утворюють раціональну базу для згаданих вище Лагранжевих областей родів.

2. Кожна раціональна база мусить (для полів першого типу) містити принаймні $n$ функцій; якщо вона містить рівно $n$ функцій, які тоді мусять бути алгебраїчно незалежними, я називаю її ``мінімальною базою''. Питання про існування мінімальної бази частково розв'язується теоремами Люрота, Кастельнуово й Енрікеса.\footnote{J. Lüroth: Beweis eines Satzes über rationale Kurven. Math. Ann. 9. 1875. --- G. Castelnuovo: Sulla razionalità delle involuzioni piane. Math. Ann. 44. 1893. --- F. Enriques: Sopra una involuzione non razionale dello spazio. Rend. Acc. Linc. vol. 21. 21. Jan. 1912.} З них випливає, що для полів від однієї невизначеної мінімальна база завжди існує: це функція Люрота цього поля, однозначно визначена з точністю до дробово-лінійного перетворення (пор. Штайніц, \S\ 24). Для полів від двох невизначених мінімальна база існує завжди і загалом лише тоді, коли допускають алгебраїчне розширення поля коефіцієнтів, тоді як для трьох і більше невизначених мінімальна база загалом взагалі не існує. Охарактеризувати спеціальні поля з мінімальною базою ще не пробували.

Тепер я далі дослідила питання про мінімальну базу для Лагранжевих областей родів. Тут воно набуває такого значення: ``Якщо Лагранжева область родів, що належить групі $G$, має мінімальну базу, то всю сукупність рівнянь із заданою групою $G$ можна раціонально побудувати за допомогою параметричного подання; і саме для кожного числового поля, яке містить область коефіцієнтів мінімальної бази -- отже, зокрема, для будь-якого числового поля, якщо ця область коефіцієнтів є полем раціональних чисел.''

Найпростіший приклад дає симетрична група; тут елементарні симетричні функції утворюють мінімальну базу з раціональними числовими коефіцієнтами. Звідси як параметричне подання одержуємо просто рівняння з невизначеними коефіцієнтами. Як від нього перейти до як завгодно багатьох рівнянь ``без афекту'' над кожним числовим полем, показує теорема Гільберта про незвідність.

Безпосередньо з теорем Люрота і Кастельнуово випливає існування мінімальної бази для всіх груп, що трапляються при рівняннях 3-го і 4-го степенів; при цьому далі виявляється, що ця мінімальна база має раціональні числові коефіцієнти. Отже, над довільним числовим полем можна раціонально побудувати сукупність рівнянь 3-го і 4-го степенів із заданою групою. Для циклічних і діедральних груп існує мінімальна база з полем $n$-тих коренів з одиниці як областю коефіцієнтів; те саме справджується для деяких інших метациклічних груп. Питання, чи мінімальна база взагалі є характерною для певних груп, мусить лишитися нерозв'язаним.

3. Щоб перейти до останнього базового поняття, я розгляну сукупність поліномів, що містяться в полі. Під ``інтегральною базою'' я розумію скінченну кількість таких поліномів, що кожний поліном поля можна подати як цілу раціональну комбінацію цієї скінченної кількості поліномів з коефіцієнтами із заданої області коефіцієнтів. Питання про існування інтегральної бази -- яке, наприклад, як спеціальний випадок містить скінченність системи інваріантів -- у дещо іншому формулюванні поставив Гільберт у своїх ``Математичних проблемах''\footnote{D. Hilbert: Mathematische Probleme. Доповідь, Париж 1900. Проблема 14. Göttinger Nachr. 1900.} як проблему відносно цілих функцій.

Наразі я можу вказати лише один клас полів, для яких інтегральна база існує. А саме, якщо поле містить систему з $n$ поліномів таку, що однорідна результанта членів найвищого степеня цих поліномів не дорівнює нулю, то інтегральна база існує; її дають коефіцієнти всіх $u$ у незвідному над полем рівнянні для лінійної форми:
\[
  u_0=u_1x_1+u_2x_2+\cdots+u_nx_n.\footnotemark
\]
\footnotetext{Це незвідне рівняння у випадку однієї невизначеної трапляється у Е. Штайніца в доведенні теореми Люрота.}
Якщо, зокрема, значення результанти дорівнює одиниці області коефіцієнтів, то забезпечена також інтегральність подання. Прикладом цього класу полів є Лагранжеві області родів, бо результанта елементарних симетричних функцій має значення $\pm 1$. Інший приклад (без інтегральності) дають усі поля, що містять лише один алгебраїчно незалежний поліном; інтегральна база тут тотожна функції Люрота найменшого підполя, яке містить усі поліноми. Так, як відомо, всі цілі раціональні проєктивні інваріанти квадратичної форми від $n$ змінних є цілими функціями дискримінанта.

Наведена умова є лише достатньою, аж ніяк не необхідною для існування інтегральної бази. Легко побудувати поля, у яких результанта будь-яких $n$ поліномів дорівнює нулю і які все ж мають інтегральну базу, задану коефіцієнтами того незвідного рівняння. Постає питання, чи, можливо, і в найзагальнішому випадку коефіцієнти того рівняння дають інтегральну базу.

\end{document}
